\documentclass[10pt,a4paper]{amsart}
\usepackage[margin=19mm]{geometry}
\usepackage{amsmath,amssymb,mathtools}
\usepackage[scr=rsfs,cal=euler]{mathalfa}
\usepackage{xfrac}
\usepackage[unicode,hidelinks,bookmarks=false]{hyperref}
\newcommand{\ie}{i.e.\ }
\newcommand{\cf}{cf.\ }
\newcommand{\resp}{respectively\ }
\newcommand{\Subsection}{\subsection}
\providecommand{\nobreakdash}{\nobreak\hbox{-}}
\setlength{\emergencystretch}{2em}
\multlinegap0pt
\usepackage[shortlabels]{enumitem}
\usepackage{amssymb}
\usepackage{mathtools}
\usepackage{tikz}
\newtheorem{lemma}{Lemma}
\newtheorem{corollary}{Corollary}
\newtheorem{theorem}{Theorem}
\newcommand{\bbC}{\mathbb C}
\newcommand{\cstar}{C*}
\newcommand{\sfC}{\mathsf C}
\DeclareMathOperator{\sfCOM}{\mathsf {Ab}}
\newcommand{\sfD}{\mathsf D}
\DeclareMathOperator{\sfvN}{\mathsf {vN}}
\title{Ulam stability for classes of nuclear \cstar-algebras}
\author{Vadim Alekseev, Ilijas Farah, Andreas Thom}
\date{Source: arXiv:2606.03757v1 (CC BY 4.0)}
\begin{document}
\maketitle

\noindent\emph{Source and licence.} Ulam stability for classes of nuclear \cstar-algebras. Version arXiv:2606.03757v1 (CC BY 4.0), distributed by arXiv under the Creative Commons Attribution 4.0 International licence, http://creativecommons.org/licenses/by/4.0/, as recorded in the arXiv licence metadata for this exact version and shown on the abstract page https://arxiv.org/abs/2606.03757v1. Full licence notice: \texttt{licenses/arxiv-2606.03757-CC-BY-4.0.txt}.\par\smallskip

\noindent\textit{Editorial scope.} Counted unit: the article's theorem T.Abelian+ (source0.tex lines 685--691). The article's complete original proof of it is delivered with it (source0.tex lines 943--1066, one \texttt{proof} environment of 124 lines, which the article places away from the statement). No mathematics is written, repaired or re-derived: the statement and the proof are the article's own lines, in the article's own notation, and no retained line is substituted or re-flowed.  Retained uncounted as the source-local context the proof needs: lines 371--383; lines 409--440; lines 721--733; lines 738--754.  Nothing was omitted from any retained span, so the package carries no omission record.  No retained line contains a citation, so the wrapper carries no bibliography and no bibliography entry.  No retained line contains a figure environment, an image-inclusion command or drawing code, so no image is delivered and the package declares no asset dependency.  Editorial formatting only, in addition: an AMS \texttt{amsart} preamble that restates the article's own declarations and macros used by the retained lines, namely the environments \texttt{lemma}, \texttt{corollary}, \texttt{theorem} on separate counters, together with the standard packages those lines need.  The retained environments print in this document as Lemma 1, Corollary 1, Theorem 1 in order of appearance, whereas the article prints the same items with the numbering of its own full document.  The complete native source of the article as distributed in the arXiv e-print for this exact version is bundled unchanged as \texttt{sources/201998/source0.tex}.
\begin{lemma}\label{L.Corrected}
		There exists a function $g(\varepsilon)\to 0$ as $\varepsilon\to 0$ with the following property. For all \cstar-algebras $A$ and $B$ and every $\varepsilon$-$*$-homomorphism $\varphi\colon A\to B$ 
 there exists a $g(\varepsilon)$-$*$-homomorphism $\psi\colon A\to B$ such that $\sup_{\|a\|\leq 1} \|\psi(a)-\varphi(a)\| \leq g(\varepsilon)$ and in addition $\psi$ satisfies the following for all $a\in A$ such that  $\|a\|\leq 1$. 
\begin{enumerate}
	\item\label{a.Corrected.norm} $\|\psi(a)\|\leq 1$. 
	\item \label{a.Corrected.*} $\psi(a^*)=\psi(a)^*$. 
	\item \label{a.Corrected.scalar} $\psi(\lambda a)=\lambda \psi(a)$ for all $\lambda\in \bbC$ such that $|\lambda|\leq 1$. 
	\item \label{a.Corrected.ctns}$\|\psi(a)\|\leq \|a\|$. 
	\item \label{a.Corrected.leq} $\psi(a)\leq \psi(b)+3\varepsilon$ for all $0\leq a\leq b\leq 1$. 
			\item \label{a.Corrected.exp} $\exp(i\psi(a))\approx_{g(\varepsilon )}\psi(\exp(i a))$ for all $a\in A_{\textrm{sa}}$ with $\|a\|\leq 1$. 
	\item \label{a.Corrected.log} $\psi(a)\approx_{g(\varepsilon)} -i \log \psi(\exp(ia))$ for all $a\in A_{\textrm{sa}}$ with $\|a\|\leq 1$.
\end{enumerate}
\end{lemma}

\begin{corollary}\label{C.UlamStableCorrected}
Suppose that $\sfC$ and $\sfD$  are classes of \cstar-algebra. Then the following are equivalent. 
\begin{enumerate}
	\item\label{1.C.Ulam.Corrected} The pair $(\sfC,\sfD)$ is Ulam stable. 
\item\label{2.C.Ulam.Corrected} There is a function $f(\varepsilon)\to 0$ as $\varepsilon\to 0$ such that for every $A\in \sfC$, $B\in \sfD$, and $\varepsilon$-$*$-homomorphism $\varphi\colon A\to B$  that for all $a,b$ in $A_1$ and all $\lambda\in \bbC_1$ satisfies
\begin{enumerate}[label=(C\arabic*)]
	\item\label{Corrected.norm} $\|\varphi(a)\|\leq 1$.
	\item \label{Corrected.*} $\varphi(a^*)=\varphi(a)^*$. 
	\item \label{Corrected.scalar} $\varphi(\lambda a)=\lambda \varphi(a)$ for all $\lambda\in \bbC$ such that $|\lambda|\leq 1$. 
	\item \label{Corrected.ctns}$\|\varphi(a)\|\leq \|a\|$. 
	\item \label{Corrected.leq} $\varphi(a)\leq \varphi(b)+3\varepsilon$ for all $0\leq a\leq b\leq 1$. 
	\item \label{Corrected.exp} $\exp(i\varphi(a))\approx_{g(\varepsilon )}\varphi(\exp(i a))$ for all $a\in A_{\textrm{sa}}$ with $\|a\|\leq 1$. 
	\item \label{Corrected.log} $\varphi(a)\approx_{g(\varepsilon)} -i \log \varphi(\exp(ia))$ for all $a\in A_{\textrm{sa}}$ with $\|a\|\leq 1$.
\end{enumerate}
%
%\begin{align}\label{eq.e-$*$-homomorphism}
%\begin{split}	\|\varphi(a+b)-\varphi(a)-\varphi(b)\| &< \varepsilon,\\
%	\|\varphi(ab)-\varphi(a)\varphi(b)\| &< \varepsilon,\\
%	\varphi(a^*)&=\varphi(a)^*\\
%	\varphi(\lambda a)&=\lambda\varphi(a)\|,\\
%\|\varphi(a)\|&\leq \|a\|,\\
%\varphi(a)&\leq \varphi(b)+\varepsilon,\quad\text{ if $0\leq a\leq b\leq 1$},\\  
% \|\exp(i\varphi(a))-\varphi(\exp(i a))\|&\leq \varepsilon \text{ if $a=a^*$},\\
%\|\varphi(a)- i \log \varphi(\exp(ia))\|&\leq \varepsilon\text{ if $a=a^*$}
%\end{split}
%\end{align}
there is a $*$-homomorphism $\psi\colon A\to B$ such that $$\sup_{\|a\|\le 1}\|\varphi(a)-\psi(a)\|\le f(\varepsilon).$$
\end{enumerate}
\end{corollary}
\begin{proof}
	By Lemma~\ref{L.Corrected}, every $\varepsilon$-$*$-homomorphism $\varphi_0$ from $A$ into $B$ can be $g(\varepsilon)$-approximated on $A_1$ by a function $\varphi$ that satisfies  conditions in \eqref{2.C.Ulam.Corrected}. Clearly $\varphi$ is uniformly approximated by a $*$-homomorphism on $A_1$ if and only if $\varphi_0$ is uniformly approximated by a $*$-homomorphism on $A_1$. 
\end{proof}

\begin{theorem}
\label{T.Closure}
The class of Ulam stable $C^*$-algebras is closed by taking:
\begin{enumerate}
\item \label{I.Ideals} norm-closed,  two-sided,  ideals,
\item \label{I.Quotients} quotients by norm-closed,  two-sided ideals,
\item \label{I.Extensions} extensions,
\item \label{I.Pullback} pullbacks where one map is surjective,
\item \label{thm:Johnson} inductive limits, and
\item \label{I.Tensor}tensor products with abelian \cstar-algebras.
\end{enumerate}
Moreover, if $\sfC$ is a class of \cstar-algebras such that the pair $(\sfC,\sfvN)$ is Ulam stable, then so is each of the classes $\sfC'$obtained from $\sfC$ by applying any of the operations listed above, with a possibly different witnessing function.
\end{theorem}

\begin{lemma}\label{L.extend.from.ideal}
Let $A$ be a \cstar-algebra, let $I$ be a closed two-sided ideal of $A$,
let $M$ be a von Neumann algebra, and let
$
\varphi\colon I\to M
$ 
be an $\varepsilon$-$*$-homomorphism.
Then there is a $3\varepsilon$-$*$-homomorphism
$
\widetilde\varphi\colon A\to M
$
such that
\[
\sup_{x\in I_1}\|\widetilde\varphi(x)-\varphi(x)\|
\leq 2\varepsilon.
\]
\end{lemma}

\begin{theorem} \label{T.Abelian+} Each of the following pairs is Ulam stable. 
	\begin{enumerate}
		\item \label{U.Abelian.vN} $(\sfCOM,\sfvN)$. 
		\item \label{U.Abelian.Abelian}
		$(\sfCOM,\sfCOM)$. 
	\end{enumerate}
\end{theorem}

\begin{proof}[Proof of Theorem~\ref{T.Closure}\eqref{I.Tensor}] 
	Fix an abelian $A$, $C\in \sfC$, and an $\varepsilon$-$*$-homomorphism $\varphi\colon A\otimes  C\to M$. Without loss of generality, we may assume that $C$ is unital.
	In the following, by  $A^+$ we denote the unitization of a \cstar-algebra $A$.  
	By Lemma~\ref{L.extend.from.ideal} we may extend $\varphi$ to $A^+\otimes C$ and assume that $A$ is unital. 
	
	By Theorem~\ref{T.Abelian+} \eqref{U.Abelian.vN}, there is a $*$-homomorphism $\Phi_A\colon A\to M$ that $f_{\sfCOM}(\varepsilon)$-approximates the restriction of $\varphi$ to the unit ball of $A$. Since $M$ is a von Neumann algebra and the unitary group of $A$ is abelian, we can define $\varphi_1\colon A\otimes C\to M$ by  
	\[
	\varphi_1(x)=\int \Phi_A(u)\varphi(x)\Phi_A(u^{-1})\, d\mu(u)
	\]
	where $\mu$ is an invariant mean on the unitary group of $A$ (considered as a discrete group). We have, for every $x\in (A\otimes C)_1$ and $u\in U(A)$ (by Corollary~\ref{C.UlamStableCorrected}   \eqref{1.C.Ulam.Corrected}, 
	 we may assume that $\varphi$ sends the norm-unit ball into the norm-unit ball)
	\begin{align*}
		\|\varphi(x)-\Phi_A(u)\varphi(x)\Phi_A(u^{-1})\|
		&\approx_{2f_{\sfCOM}(\varepsilon)} 
		\|\varphi(x)-\varphi(u)\varphi(x)\varphi(u^{-1})\|\\
		&\approx_{2\varepsilon}
		\|\varphi(x)-\varphi(uxu^{-1})\|=0. 
	\end{align*}
	Since $\varphi_1$ is $2(f_{\sfCOM}(\varepsilon)+\varepsilon)$-approximated by an $\varepsilon$-$*$-homomorphism $\varphi$, it is a $\delta$-$*$-homomorphism, with 
	\[
	\delta=3(f_{\sfCOM}(\varepsilon)+\varepsilon)+\varepsilon. 
	\]
	We gained centrality at the cost of increasing the error, i.e.
	$$\Phi_A(a)\varphi_1(x) = \varphi_1(x)\Phi_A(a) \quad (x \in A \otimes C, \quad a \in A).$$
	
	Note that the range of $\varphi_1$ is included in the commutant
	$\Phi_A[A]'$, which is again a von Neumann algebra.
	By Ulam stability of $\sfC$ with respect to von Neumann algebras,
	the restriction of $\varphi_1$ to $1\otimes C$ can be
	$f_{\sfC}(\delta)$-approximated by a $*$-homomorphism
	$
		\Phi_C\colon C\to \Phi_A[A]'.
	$
	Since the ranges of $\Phi_A$ and $\Phi_C$ commute, we obtain a
	$*$-homomorphism
	$
		\Phi=\Phi_A\otimes \Phi_C\colon A\otimes C\to M.
	$
	We claim that $\Phi$ approximates $\varphi_1$ uniformly on the unit ball.

	Let $X$ be the spectrum of $A$, so that $A=C(X)$.
	Set
	\[
		N:=W^*\bigl(\varphi_1[A\otimes C],\,\Phi_A[A],\,\Phi_C[C]\bigr)
		\subseteq M.
	\]
	Since $\varphi_1[A\otimes C]\subseteq \Phi_A[A]'$ and
	$\Phi_C[C]\subseteq \Phi_A[A]'$, the algebra $\Phi_A[A]$ is contained
	in the centre of $N$. Hence $N$ is a $C(X)$-algebra. For $x\in X$, let
	$
		N_x:=N/C_0(X,x)N
	$
	and denote the quotient map by $q_x\colon N\to N_x$.
	For $z\in N$ we have
	\[
		\|z\|=\sup_{x\in X}\|q_x(z)\|.
	\]

	Note that$\Phi_A(a)\varphi_1(x) \approx_{2\delta} \varphi_1(ax)$ for all $x \in A \otimes C$ and $a \in A$.
	We first record a consequence of this approximate compatibility with the $C(X)$-module structure.
	If $d\in (A\otimes C)_1$ satisfies $d(x)=0$, then
	$
		\|q_x(\varphi_1(d))\|\leq 4\delta.
	$
	Indeed, given $\eta>0$, choose $a\in A_+$ with
	$0\leq a\leq 1$, $a(x)=0$, and $\|d-ad\|<\eta$.
	Using approximate additivity and approximate homogeneity of
	$\varphi_1$, we have
	$
		\|\varphi_1(d)-\varphi_1(ad)\|
		\leq \eta+2\delta.
	$
	On the other hand,
	$
		q_x(\varphi_1(ad))
		\approx_{2\delta}q_x(\Phi_A(a)\varphi_1(d))=0,
	$
	because $a(x)=0$. Letting $\eta\to 0$ gives the claim.

	Now fix $b\in (A\otimes C)_1$ and $x\in X$, and put
	$
		c_x:=1\otimes b(x)\in A\otimes C.
	$
	Then $e_x:=\frac12(b-c_x)$ belongs to $(A\otimes C)_1$ and satisfies
	$e_x(x)=0$. By the preceding claim,
	$
		\|q_x(\varphi_1(e_x))\|\leq 4\delta.
	$
	Using approximate additivity and approximate homogeneity once more,
	\begin{align*}
		\|q_x(\varphi_1(b)-\varphi_1(c_x))\|
		&\leq \|q_x(\varphi_1(b-c_x))\|+2\delta\\
		&=\|q_x(\varphi_1(2e_x))\|+2\delta\\
		&\leq 2\|q_x(\varphi_1(e_x))\|+3\delta\\
		&\leq 10\delta.
	\end{align*}
	Since $\Phi$ is a $C(X)$-homomorphism,
	$
		q_x(\Phi(b))=q_x(\Phi(c_x))
		=q_x(\Phi_C(b(x))).
	$
	Moreover, by the choice of $\Phi_C$,
	$
		\|\varphi_1(c_x)-\Phi_C(b(x))\|
		\leq f_{\sfC}(\delta).
	$
	Therefore,
	\begin{align*}
		\|q_x(\varphi_1(b)-\Phi(b))\|
		&\leq
		\|q_x(\varphi_1(b)-\varphi_1(c_x))\|
		+\|q_x(\varphi_1(c_x)-\Phi_C(b(x)))\|\\
		&\leq 10\delta+f_{\sfC}(\delta).
	\end{align*}
	Taking the supremum over $x\in X$, we obtain
	\[
		\|\varphi_1(b)-\Phi(b)\|
		\leq 10\delta+f_{\sfC}(\delta)
		\qquad (b\in (A\otimes C)_1).
	\]
	Since $\varphi_1$ uniformly approximates $\varphi$ on the unit ball
	by a quantity tending to $0$ with $\varepsilon$, the same is true of
	$\Phi$. This finishes the proof.
\end{proof}

\end{document}
